\section{Síntesis y ampliación}

Esta entrevista se puede condensar en una tesis general: \textbf{el valor a largo plazo de un Agente de propósito general proviene de la capacidad sistémica de «completar de forma estable tareas complejas de cola larga», no del dividendo puntual de una versión de modelo.}

\subsection{Tabla comparativa de conclusiones centrales}

\begin{center}
\begin{tabular}{p{3.4cm}p{5.1cm}p{5.3cm}}
\toprule
Tema & Conclusión de la entrevista & Acción ejecutable \\
\midrule
Posicionamiento del producto & Primero general, después vertical; la cola larga es la fuente de crecimiento incremental & Recopilar primero la distribución real de tareas, y luego decidir la profundización vertical \\
Elección de arquitectura & Agente único + entorno aislado + centro de codificación & Unificar el espacio de acciones, reducir la complejidad innecesaria de herramientas \\
Ciclo de evaluación & Evaluación automatizada + revisión humana + retroalimentación del usuario & Construir un repositorio de muestras de fallo y un sistema de reproducción \\
Mecanismo organizativo & GPA + prioridad a la acción + no hacer listas & Convergencia de objetivos semanal, retrospectiva quincenal que elimina direcciones poco eficientes \\
Estrategia competitiva & Las funciones se homogeneizan; la eficiencia del sistema es el foso real & Invertir en velocidad de iteración y entrega confiable \\
\bottomrule
\end{tabular}
\end{center}

\begin{importantbox}{Tres prioridades para la implementación en equipo}
\begin{enumerate}
    \item Fijar la «tasa de éxito de la tarea» como el indicador principal, no un benchmark aislado
    \item Conservar una traza reproducible de cada fallo, para reducir el costo de la siguiente corrección
    \item Depurar periódicamente funciones y deuda de herramientas, para mantener una complejidad del sistema controlable
\end{enumerate}
\end{importantbox}

\subsection{Lecturas adicionales}

\begin{itemize}
    \item Materiales oficiales de lanzamiento y demostraciones públicas de Manus
    \item Word2Vec (Mikolov et al., 2013)
    \item BERT (Devlin et al., 2018) y la arquitectura transformer
    \item Reseñas relacionadas con Open Information Extraction
    \item Benchmarks públicos y artículos sobre Agent Evaluation y Tool Use Reliability
\end{itemize}

%% --- Fin del contenido --- %%

\end{document}
